\documentclass[article,36pt,extrafontsizes,oneside,openany,oldfontcommands]{memoir}
\usepackage{graphicx}
    \graphicspath{{figure/}}
\usepackage{amsmath, amsfonts, amssymb, amsthm}
\usepackage{mathtools}
\usepackage[T1]{fontenc}
\usepackage[utf8]{inputenc}
\usepackage{lmodern}
\usepackage{bm}
\RequirePackage{lipsum}
\RequirePackage{blindtext}
\RequirePackage[svgnames,table]{xcolor}
\RequirePackage{tikz}
\RequirePackage[framemethod=tikz]{mdframed}
\RequirePackage{color}
\RequirePackage{geometry}
\RequirePackage{adjmulticol}
\RequirePackage[skins,most,listings,skins]{tcolorbox}

%For table formatting
\RequirePackage{longtable}
\RequirePackage{array}
\RequirePackage{multirow}
\RequirePackage{wrapfig}
\RequirePackage{float}
\RequirePackage{colortbl}
\RequirePackage{pdflscape}
\RequirePackage{pagecolor}
\RequirePackage{tabu}
\RequirePackage{threeparttablex}
\RequirePackage[normalem]{ulem}
\RequirePackage{makecell}
\usepackage{xmpmulti, booktabs, multicol}
\usepackage{threeparttable}
\usepackage{dcolumn}
\newcolumntype{.}[1]{D{.}{.}{#1}}
\newcolumntype{,}[1]{D{,}{,}{#1}}
\usepackage{siunitx}
\sisetup{%
  mode = math,
  detect-all,  
  output-decimal-marker = {.},
  group-separator = {,},
  math-rm = \rmnumeric,
  text-rm = \rmnumeric,
}

% Color definitions
\usepackage{color}
    \definecolor{MyBrown}{HTML}{3D3D00}
    \definecolor{MyBlue}{HTML}{003A75}  
    \definecolor{MyRed}{HTML}{6F1A2F}
    \definecolor{MyGreen}{HTML}{004847}
\usepackage{hyperref}
    \hypersetup{bookmarks = true,
                bookmarksnumbered = true,
                colorlinks = true,
                linkcolor = MyRed,
                citecolor = MyBlue,
                filecolor = MyBrown,
                urlcolor = MyGreen}

\RequirePackage{float}
\floatplacement{figure}{H}
\floatplacement{table}{H}

%%%%%%%%% COLOURS %%%%%%%%
\definecolor{titleboxbgcol}{HTML}{D8E0FF}
\definecolor{titleboxbordercol}{HTML}{205DA4}
\definecolor{columnlinecol}{HTML}{008080}
\definecolor{bodybgcol}{HTML}{ffffff}
\definecolor{sectitlebgcol}{HTML}{D8E0FF}
\definecolor{sectitlebordercol}{HTML}{D8E0FF}
% Text Colours
\definecolor{titletextcol}{HTML}{1c518f}
\definecolor{authortextcol}{HTML}{000000}
\definecolor{affiliationtextcol}{HTML}{000000}
\definecolor{sectitletextcol}{HTML}{1c518f}
\definecolor{bodytextcol}{HTML}{000000}
\definecolor{footnotetextcol}{HTML}{000000}
\definecolor{citecol}{HTML}{CC0000}
\definecolor{urlcol}{HTML}{008080}
\definecolor{linkcol}{HTML}{008080}

%Memoir spacing options
%spacing between figure/table and caption
\setlength{\abovecaptionskip}{0.4in}
\setlength{\belowcaptionskip}{0.2in}
\captionnamefont{\large}
\captiontitlefont{\large}

%define column options
\setlength{\columnseprule}{0pt}
\def\columnseprulecolor{\color{columnlinecol}}

%define section title features
\setsubsubsecheadstyle{\small\color{sectitletextcol}\textbf}
\setsecnumformat{}
\def\sectionmark#1{\markboth{#1}{#1}}

%%%%%%%%%%%% TCOLORBOXES %%%%%%%%%%%%%%%%%%%%
%Title Box
\newtcolorbox{topbox}{
enhanced,
colback=titleboxbgcol,
colframe=titleboxbordercol,
halign=center,
boxrule=0cm,
sharp corners=all,
}

%Body Section Title Box
\newtcolorbox{myboxstuff}[1][]{
code={\parindent=0em},
colframe=sectitlebordercol,
nobeforeafter,
left skip=0pt,
valign=center,
halign=center,
fontupper=\Large\bfseries,
colupper=sectitletextcol,
boxrule=2mm,
colback=sectitlebgcol,
sharp corners=south, #1}
\newcommand{\mybox}[1]{%
\begin{myboxstuff}
\strut #1
\end{myboxstuff}%
}
\makeheadstyles{MyBox}{
    \setsecheadstyle{\mybox}
}
\headstyles{MyBox}\makepagestyle{MyBox}

\thispagestyle{empty}

%Remove section numbering
\counterwithout{section}{chapter}
\makechapterstyle{mydefault}{
\addtocounter{secnumdepth}{2}
\setsecheadstyle{\mybox}
\setsubsecheadstyle{\itshape}
\setsubsubsecheadstyle{\itshape}
}

%set the chapterstyle
\chapterstyle{mydefault}

%define column spacing
\setlength\columnsep{0.5in}

%spacing params
\setlength\parindent{0em}
\setlength\parskip{0em}
\setlength\hangparas{0em}

%spacing after section head title
\setaftersecskip{0em}
\setbeforesecskip{1.5em}
\setlength\textfloatsep{0in}
\setlength\floatsep{0.1in}
\setlength\intextsep{0.5in}
\makeatletter
\g@addto@macro{\normalsize}{%
    } 
\makeatother

\setstocksize{36in}{48in}
\settrimmedsize{\stockheight}{\stockwidth}{*}
\settypeblocksize{36in}{48in}{*}
\setlrmargins{*}{*}{1}
\setulmarginsandblock{2.5cm}{0cm}{*}
\setmarginnotes{0em}{0cm}{0cm}
\setlength{\footskip}{0cm}
\setlength{\footnotesep}{0cm}
\setlength{\headheight}{0pt}
\setlength{\headsep}{0pt}
\setlength{\trimtop}{0pt}
\setlength{\trimedge}{0pt}
\setlength{\uppermargin}{0pt}
\checkandfixthelayout
\usepackage{enumitem}
\usepackage{bbm}
\usepackage{dsfont}
\def\one{\mathrm{1}\hspace{-12pt}\mathrm{1}}
\setlist[enumerate]{topsep=0pt,itemsep=-1ex,partopsep=1ex,parsep=1ex,leftmargin=2ex}
\setlist[itemize]{topsep=0pt,itemsep=-1ex,partopsep=1ex,parsep=1ex,leftmargin=2ex}

%Footnote styling
\RequirePackage{footmisc}
\def\footnotelayout{\centering\color{footnotetextcol}}

% Code highlighting for markdown
\usepackage{color}
\usepackage{fancyvrb}
\newcommand{\VerbBar}{|}
\newcommand{\VERB}{\Verb[commandchars=\\\{\}]}
\DefineVerbatimEnvironment{Highlighting}{Verbatim}{commandchars=\\\{\}}
\usepackage{framed}
\definecolor{shadecolor}{RGB}{248,248,248}
\newenvironment{Shaded}{\begin{snugshade}}{\end{snugshade}}

% Suppress warnings about fonts
\usepackage{fancyhdr}

% Font family - using standard LaTeX fonts
\renewcommand{\familydefault}{\rmdefault}  % Use serif (Roman) font

% define the BODYBGCOL
\newpagecolor{bodybgcol}

%sets footnote styling
\renewcommand\footnoterule{}

%-------------- Begin Document -------------------%
\begin{document}

%-------------- Title Box Start ------------------%
\begin{topbox}
  \color{titletextcol}
  \vspace{1.2in}
  \bfseries
  \fontsize{115pt}{115pt}\selectfont
  {Evaluating In-Context Learning Capabilities of Transformers on Binary Classification}\\[0.75in]
  \color{authortextcol} 
  \mdseries
  \LARGE{Rushil Patel, Banga Bansal $\cdot$ UC San Diego $\cdot$ DSC 180 $\cdot$ Spring 2025} 
  \vspace{0.75in}
\end{topbox}

%--------------- Title Box End -------------------%
%----------------- Body Start --------------------%

\fontsize{45pt}{55pt}\selectfont

\begin{adjmulticols*}{3}{0.5in}{0.5in}
\color{bodytextcol}

\section{Research Overview}

\begin{adjustwidth}{0.1in}{0.1in}
In-context learning (ICL) enables large language models and transformers to perform novel tasks by inferring patterns from in-context examples without parameter updates. This project investigates ICL performance on synthetic binary classification tasks from Gaussian mixtures, addressing three key research questions:

\vspace{0.2in}
\begin{itemize}[topsep=0pt,itemsep=0ex,partopsep=1ex,parsep=1ex]
\item \textbf{RQ1:} How does in-context test accuracy scale with dimension $d$, batch size $B$, and sequence length $N$?

\item \textbf{RQ2:} Under what conditions does the model exhibit \emph{benign overfitting}, memorizing noisy labels while maintaining high test accuracy?

\item \textbf{RQ3:} Are phenomena from linear attention transformers preserved with full transformer architectures (encoder-only, decoder-only)?
\end{itemize}

\end{adjustwidth}

\section{Data Generation: Gaussian Mixture Model}

\begin{adjustwidth}{0.1in}{0.1in}
For each classification task $\tau$, we generate data as follows:

\vspace{0.1in}
\begin{itemize}[topsep=0pt,itemsep=0ex,partopsep=1ex,parsep=1ex]
\item Sample a task-specific mean: $\mu_\tau \sim \text{Unif}(R \cdot S^{d-1})$
\item Sample labels: $y_{\tau,i} \in \{-1, +1\}$ uniformly
\item Sample noise: $z_{\tau,i} \sim \mathcal{N}(0, I_d)$
\item Form input: $x_{\tau,i} = y_{\tau,i} \mu_\tau + z_{\tau,i}$
\end{itemize}

\vspace{0.15in}
Each task consists of:
\begin{itemize}[topsep=0pt,itemsep=0ex,partopsep=1ex,parsep=1ex]
\item $N$ in-context labeled examples $(x_1, y_1), \ldots, (x_N, y_N)$
\item Signal strength controlled by $R$ (class separability)
\item Label noise via probability $p$ of flipping
\end{itemize}

\end{adjustwidth}

\columnbreak

\section{In-Context Learning Mechanism}

\begin{adjustwidth}{0.1in}{0.1in}
The model receives a sequence of labeled examples and must predict:
$$\hat{y}_{N+1} = \text{sign}(\text{model}(\{(x_i, y_i)\}_{i=1}^{N}, x_{N+1}))$$

\vspace{0.2in}
\textbf{Linear Transformer:} One-layer linear attention model with:
\begin{align*}
W &: \text{learnable weight matrix} \\
\text{Prediction} &= \text{sign}\left(x_{N+1}^T W \frac{\sum_i y_i x_i}{\sum_i y_i^2}\right)
\end{align*}

\vspace{0.2in}
\textbf{Benign Overfitting:} Occurs when the model:
\begin{itemize}[topsep=0pt,itemsep=0ex,partopsep=1ex,parsep=1ex]
\item Memorizes noisy in-context labels (high training accuracy)
\item Yet generalizes well to clean test data (high test accuracy)
\end{itemize}

This paradoxical behavior mirrors modern large language models.

\end{adjustwidth}

\section{RQ1: Scaling with Dimension, Batch Size, and Sequence Length}

\begin{adjustwidth}{0.1in}{0.1in}
\textbf{Experimental Setup:}
\begin{itemize}[topsep=0pt,itemsep=0ex,partopsep=1ex,parsep=1ex]
\item Vary dimension: $d \in \{10, 50, 100, 500, 1000\}$
\item Vary sequence length: $N \in \{5, 10, 20, 40, 80\}$
\item Vary batch size: $B \in \{50, 100, 250, 500, 1000, 2000\}$
\item Signal strength: $R = \sqrt{d}^{\alpha}$ for various $\alpha$
\end{itemize}

\vspace{0.15in}
\textbf{Expected Phenomena:}
\begin{itemize}[topsep=0pt,itemsep=0ex,partopsep=1ex,parsep=1ex]
\item Accuracy improves with sequence length $N$ (more examples)
\item Accuracy depends on $R/\sqrt{d}$ (signal-to-noise ratio)
\item Double descent behavior in scaling regime
\end{itemize}

\textbf{Results location:} \texttt{results/rq1/rq1\_linear\_scaling/}

\end{adjustwidth}

\columnbreak

\begin{minipage}{.65\linewidth}

\section{RQ2: Benign Overfitting Phenomenon}

\begin{adjustwidth}{0.1in}{0.1in}
\textbf{Key Investigation:} Does the model memorize noisy labels while generalizing?

\vspace{0.15in}
We measure accuracy on:
\begin{itemize}[topsep=0pt,itemsep=0ex,partopsep=1ex,parsep=1ex]
\item \textbf{Training (noisy):} Accuracy on flipped labels ($p \in \{0, 0.2, 0.4\}$)
\item \textbf{Test (clean):} Accuracy on unseen clean data
\end{itemize}

\vspace{0.15in}
\textbf{Hypothesis:} In high-dimension, low-noise regime:
\begin{itemize}[topsep=0pt,itemsep=0ex,partopsep=1ex,parsep=1ex]
\item Training accuracy may approach or reach 100\% even with noise
\item Test accuracy remains close to 50\% + margin from signal
\end{itemize}

This suggests the model learns task-invariant structure despite noise.

\vspace{0.15in}
\textbf{Evaluation Metrics:}
\begin{itemize}[topsep=0pt,itemsep=0ex,partopsep=1ex,parsep=1ex]
\item Train-test gap as function of label flip probability
\item Phase transitions in accuracy curves
\end{itemize}

\end{adjustwidth}

\end{minipage}

\section{RQ3: Full Transformer Architectures}

\begin{adjustwidth}{0.1in}{0.1in}
\textbf{Extensions Beyond Linear Attention:}

\vspace{0.15in}
We test whether ICL phenomena generalize to:
\begin{itemize}[topsep=0pt,itemsep=0ex,partopsep=1ex,parsep=1ex]
\item \textbf{Decoder-only:} Standard GPT-like architecture
\item \textbf{Encoder-only:} BERT-like architecture  
\item \textbf{Hybrid:} Encoder-decoder similar to T5
\end{itemize}

\vspace{0.15in}
\textbf{Comparisons:}
\begin{itemize}[topsep=0pt,itemsep=0ex,partopsep=1ex,parsep=1ex]
\item How does attention (multi-head vs.\ single) affect ICL?
\item Role of positional encodings in in-context learning
\item Computational efficiency trade-offs
\end{itemize}

\textbf{Commercial LLM Evaluation:} We also evaluate Gemini, Claude, and GPT on the same tasks via API to understand how real LLMs perform on synthetic ICL tasks.

\end{adjustwidth}

\section{Methodology}

\begin{adjustwidth}{0.1in}{0.1in}
\textbf{Training Setup:}
\begin{itemize}[topsep=0pt,itemsep=0ex,partopsep=1ex,parsep=1ex]
\item Optimizer: AdamW with learning rate $5 \times 10^{-4}$
\item Loss: Binary cross-entropy
\item Training steps: 300 per task
\item Validation: Separate clean validation set
\end{itemize}

\vspace{0.15in}
\textbf{Evaluation:}
\begin{itemize}[topsep=0pt,itemsep=0ex,partopsep=1ex,parsep=1ex]
\item Accuracy metric: Percentage of correct predictions
\item Multiple seeds for reproducibility
\item Aggregated results by configuration
\end{itemize}

\textbf{Implementation:} Python with PyTorch and Transformers library. Models trained on CPU and GPU. Results saved as JSONL for analysis and visualization.

\end{adjustwidth}

\vspace{1.5in}

\section{Expected Contributions}

\begin{adjustwidth}{0.1in}{0.1in}
\begin{enumerate}[topsep=0pt,itemsep=0ex,partopsep=1ex,parsep=1ex]
\item Comprehensive scaling laws for ICL with transformers
\item Empirical characterization of benign overfitting phenomenon
\item Comparison of architectures for in-context learning
\item Evaluation of commercial LLMs on controlled ICL tasks
\item Open-source codebase for reproducing results
\end{enumerate}
\end{adjustwidth}

% \bibliographystyle{bib/econ_aer}
% \bibliography{bib/All_Papers}

\end{adjmulticols*}
%------------------ Body End ---------------------%

\end{document}
